% Hexagon note, 24 Nov 2026. CC0 1.0.
% Mathematics: GPT-6 Astra (ChatGPT Pro mode). Text, proofs as written here, and the enclosure program:
% Claude Opus 5.5 at effort high (checker and writer).
\documentclass[11pt,a4paper]{article}
\usepackage[T1]{fontenc}
\usepackage{lmodern}
\usepackage{microtype}
\usepackage[margin=2.7cm]{geometry}
\usepackage{amsmath,amssymb,amsthm}
\usepackage{booktabs}
\PassOptionsToPackage{hyphens}{url}
\usepackage[hidelinks]{hyperref}
\emergencystretch=2em

\theoremstyle{plain}
\newtheorem{theorem}{Theorem}
\newtheorem{corollary}[theorem]{Corollary}
\newtheorem{lemma}[theorem]{Lemma}
\newtheorem{proposition}[theorem]{Proposition}
\theoremstyle{remark}
\newtheorem{remark}[theorem]{Remark}

\DeclareMathOperator{\tr}{tr}
\DeclareMathOperator{\ind}{ind}
\newcommand{\R}{\mathbb R}
\newcommand{\C}{\mathbb C}
\newcommand{\Hs}{\mathcal H}
\newcommand{\workbenchurl}{https://github.com/KokunoYumeto/zeta-and-l-functions-workbench/tree/main/work/WB-0009/v1}
\newcommand{\recordurl}{https://github.com/KokunoYumeto/mathematics-workbenches/blob/main/works/WB-0009.md}

\title{A finite arithmetic test for simplicity of the Riemann zeros,\\ after Suzuki}
\author{}
\date{24 November 2026 \quad CC0 1.0}

\begin{document}
\maketitle

\begin{abstract}
\noindent
M.~Suzuki proved that, under the Riemann hypothesis, the compressions of the inverse of a self-adjoint
differentiation operator to windows of Weil's quadratic form converge in Hilbert--Schmidt norm, and that their
eigenvalues recover every distinct zero ordinate of the Riemann zeta function; he notes that the multiplicities are
not recovered. The Riemann hypothesis is now a theorem of L.~Alp\"oge. We show that the multiplicities are recovered
by even trace powers. Let $S_k=2\sum_{\gamma>0}m_\gamma\gamma^{-2k}$ be the Taylor moments of $\log\xi(1/2+iz)$,
which count multiplicity, and let $D_{N,k}$ be $S_k$ minus the trace of the $2k$-th power of a finite compression
whose matrix entries come from the explicit formula with finitely many prime powers. Then
$D_{N,k}\ge 2E(T)/T^{2k}$, where $E(T)$ is the number of repeated zeros up to height $T$ counted with their excess
multiplicity, and $\sup_{N,k}(2/D_{N,k})^{1/(2k)}$ equals the height of the first multiple zero. So the zeros up to a
height $T$ are simple if and only if one finite arithmetic inequality holds, and all zeros are simple if and only if
these heights are unbounded. An equivalent integer form counts the positive eigenvalues of one Hermitian matrix. A
computation in ball arithmetic, with prime powers below $e^{12}$ and no zero ordinates as input, certifies that every
zero with $0<\Im\rho\le 97.632$ is simple. The mathematics was produced by GPT-6 Astra (ChatGPT Pro mode); Claude
Opus~5.5 at effort high checked every proof, wrote this text, and recomputed the certificate with an independently
written program. No human has checked this note.
\end{abstract}

\section{The question and the theorem}

\paragraph{The setting.}
Let $\xi(s)=\frac12 s(s-1)\pi^{-s/2}\Gamma(s/2)\zeta(s)$ and $\Xi(z)=\xi(\tfrac12+iz)$. Alp\"oge's theorem
\cite{alpoge,lean} gives the Riemann hypothesis: every zero of $\Xi$ is real. Let $\Gamma\subset\R$ be the set of
distinct zeros of $\Xi$ and $m_\gamma$ the multiplicity of $\gamma\in\Gamma$. Since $\Xi$ is even and
$\Xi(0)=\xi(1/2)\ne0$, $\Gamma=-\Gamma$, $m_{-\gamma}=m_\gamma$, $0\notin\Gamma$, and, because $\Xi$ has order one,
$\sigma:=\sum_{\gamma\in\Gamma}m_\gamma\gamma^{-2}<\infty$ \cite[(2.2)]{suzuki}. Whether every $m_\gamma$ equals
one is the simplicity question for the Riemann zeros. It is open. Write
\[
 E(T)=\sum_{0<\gamma\le T}(m_\gamma-1),\qquad
 n(T)=\sum_{0<\gamma\le T}m_\gamma,\qquad
 \gamma_*=\min\{\gamma>0:m_\gamma\ge2\}\in(0,\infty],
\]
the multiplicity excess up to height $T$, the counted number of zeros up to $T$, and the first multiple ordinate
($\gamma_*=\infty$ if there is none).

\paragraph{Suzuki's compressions.}
Weil's quadratic form, $Q(u,v)=\sum_{\gamma\in\Gamma}m_\gamma\widehat u(\gamma)\overline{\widehat v(\gamma)}$ with
$\widehat u(t)=\int u(x)e^{itx}\,dx$, is positive, and its completion is realized by the sampling map
$u\mapsto(\widehat u(\gamma))_{\gamma\in\Gamma}$ in $\Hs=\ell^2(\Gamma,m)$, with norm
$\|c\|^2=\sum_\gamma m_\gamma|c_\gamma|^2$. On $\Hs$, let $C$ be multiplication by $1/\gamma$. Suzuki
\cite[§2.3 and (4.1)]{suzuki} shows that $C$ is a Hilbert--Schmidt operator with one eigenvector for each distinct
ordinate, and that its compressions to the subspaces of test functions supported in $[-a,a]$ converge to $C$ in
Hilbert--Schmidt norm; their positive eigenvalues decrease to $1/\gamma_1>1/\gamma_2>\cdots$ over the distinct
ordinates \cite[Theorem~1.1]{suzuki}. He notes that ``the multiplicities of the zeros are not recovered''
\cite[§1, after Theorem~1.2]{suzuki}.

We use finite-dimensional windows. For $N\ge1$, let $h=2^{-N}$ and let $V_N$ be the space of continuous functions on
$\R$, linear on each interval $[jh,(j+1)h]$, vanishing outside $(-N,N)$, with $\int f=0$. It has dimension
$d_N=2N2^N-2$. Let $P_N$ be the orthogonal projection of $\Hs$ onto the image of $V_N$ under the sampling map, and
$C_N=P_NCP_N$. For $k\ge1$ put
\[
 S_k=-2k\,[z^{2k}]\log\frac{\Xi(z)}{\Xi(0)}=2\sum_{\gamma>0}\frac{m_\gamma}{\gamma^{2k}},\qquad
 D_{N,k}=S_k-\tr C_N^{2k},\qquad H_{N,k}=\Bigl(\frac2{D_{N,k}}\Bigr)^{1/(2k)} .
\]
The middle equality is the Hadamard product (Section~\ref{sec:proof1}). Both terms of $D_{N,k}$ are computable from
arithmetic data: $S_k$ from the Taylor series of $\xi$ at $1/2$, and $\tr C_N^{2k}$ from two $d_N\times d_N$
matrices whose entries are given by the explicit formula with the prime powers below $e^{2N}$
(Section~\ref{sec:arith}).

\begin{theorem}\label{thm:main}
Assume the Riemann hypothesis (a theorem, \cite{alpoge}). For all $N,k\ge1$ and $T>0$,
\[
 D_{N,k}>2\sum_{\gamma>0}\frac{m_\gamma-1}{\gamma^{2k}}\ \ge\ \frac{2E(T)}{T^{2k}} .
\]
For fixed $k$, $D_{N,k}$ decreases in $N$ to $2\sum_{\gamma>0}(m_\gamma-1)\gamma^{-2k}$, and
\[
 \sup_{N,k\ge1}H_{N,k}=\gamma_* .
\]
Consequently, for each $T>0$, every zero with $0<\Im\rho\le T$ is simple if and only if $D_{N,k}<2/T^{2k}$ for some
$N,k\ge1$; and every nontrivial zero of $\zeta$ is simple if and only if $\sup_{N,k}H_{N,k}=\infty$.
\end{theorem}

The test has an integer form. Let $b_j=t_j-t_0$ ($0<|j|<N2^N$), where $t_j$ is the hat function of height one centred
at $jh$ with support $[(j-1)h,(j+1)h]$. These $b_j$ form a basis of $V_N$: each has mean zero, since every $t_j$
has integral $h$; they are linearly independent, since $\sum_jc_jb_j$ takes the value $c_j$ at $jh$; and a function
$f\in V_N$ equals $\sum_{|j|<N2^N}f(jh)t_j$ with $\sum_jf(jh)=h^{-1}\int f=0$, hence equals $\sum_{j\ne0}f(jh)b_j$
(this argument is that of the independent re-proof acknowledged in the note on method). Let $G_N$ and $H_N$ be the Hermitian
matrices $(G_N)_{ij}=Q(b_j,b_i)$ and $(H_N)_{ij}=\langle C\,\widehat b_j,\widehat b_i\rangle_{\Hs}$, where
$\widehat b$ denotes the sampling vector, and let $\ind_+$ denote the number of positive eigenvalues.

\begin{theorem}\label{thm:integer}
Assume the Riemann hypothesis. Let $T>0$ with $\Xi(T)\ne0$, and $\iota_N(T)=\ind_+(TH_N-G_N)$. Then $\iota_N(T)$ is
nondecreasing in $N$, at most the number of distinct ordinates in $(0,T]$, and equal to it for all large $N$. Hence
the zeros with $0<\Im\rho\le T$ are all simple if and only if $\iota_N(T)\ge n(T)$ for some $N$.
\end{theorem}

\begin{theorem}\label{thm:cert}
Every zero $\rho$ of $\zeta$ with $0<\Im\rho\le 97.632$ is simple. More precisely, $D_{6,3}<2/(97.632)^6$, and the
five rows of Table~\ref{tab:cert} with $N\le5$ satisfy the strict test at the heights shown.
\end{theorem}

The certificate uses $15040$ prime powers, the largest $162751<e^{12}$, and the matrices have dimension $766$. It
was computed twice, by two independently written programs in Arb ball arithmetic (Section~\ref{sec:cert}).

\begin{proposition}\label{prop:const}
Every zero with ordinate $\gamma\ge10$ has multiplicity $m_\gamma<\frac34\log\gamma$. The first multiple ordinate
satisfies $\gamma_*>97.632$, so $\gamma_*^2/2>4766.003712$, and every multiple zero has
$\gamma>\max(97.632,\,e^{4m_\gamma/3})$.
\end{proposition}

\paragraph{What this means.}
Theorem~\ref{thm:main} turns the simplicity question into a statement about finitely many prime powers at a time: the
multiplicity excess is exactly what separates the moments $S_k$, which see each zero with its multiplicity, from the
traces of Suzuki's compressions, which see each distinct zero once. At every fixed height the test terminates exactly
when the zeros there are simple. Simplicity of all zeros is the statement that these finite arithmetic tests succeed
at every height; the theorem gives no estimate that would decide it. The limit of the defect is the Taylor data of
$\log(\Xi(z)/(\Xi(0)\mathsf P(z)))$, where $\mathsf P(z)=\prod_{\gamma>0}(1-z^2/\gamma^2)$ is Suzuki's product over
distinct ordinates \cite[Theorem~1.2]{suzuki}; that quotient is constant exactly when all zeros are simple.

\paragraph{Why it matters.}
Simplicity of the zeros is the next question about their location after the Riemann hypothesis. It enters the
explicit formula with multiplicities, the size of $1/\zeta'(\rho)$, and the spectral interpretations of the zeros:
a self-adjoint operator whose spectrum is the set of zeros sees each one with the multiplicity of its eigenspace. Here
the comparison between the counted and the distinct zeros becomes a positive defect of explicit finite matrices, and
the defect is computable from primes.

\paragraph{The credit.}
The Hilbert space, the operator $C$, its Hilbert--Schmidt property, and the convergence of compressions are
Suzuki's \cite[§2.3, (4.1), Theorem~3.3]{suzuki}; we apply his argument to the dyadic windows $V_N$, for which we
prove the density it needs (Lemma~\ref{lem:dense}). The explicit formula is Weil's, in Bombieri's normalization as
recorded by Connes and Consani \cite[Appendix, Explicit formula]{cc-weil}. Alp\"oge's theorem \cite{alpoge}
supplies the real zeros. The principle that a finite restriction of Weil's form sees distinct zeros through its
rank or positive index and counted zeros through its trace is Alp\"oge and Furman's \cite[Proposition~4.1,
p.~6]{alpoge-furman}, and the pull-back of inertia used in Theorem~\ref{thm:integer} is their
\cite[Lemma~3.1, p.~5]{alpoge-furman}. What is new here is the exact form of that comparison for Suzuki's
compressions: the trace-power defect and its supremum (Theorem~\ref{thm:main}), the equality of the positive index
with the number of distinct ordinates at each height (Theorem~\ref{thm:integer}), the certificate
(Theorem~\ref{thm:cert}), Propositions~\ref{prop:const} and~\ref{prop:size}, and the heat-flow limit of
Proposition~\ref{prop:heat}, whose proof uses the stability criterion of Borcea and Br\"and\'en \cite{bb}.

\paragraph{Context.}
Alp\"oge and Furman \cite{alpoge-furman} proved, unconditionally, that more than two thirds of the zeros are simple
and on the critical line and more than five sixths are distinct \cite[Theorem~A, p.~1]{alpoge-furman}. They
restrict Weil's form to about $N(T,2T)$ modulated copies of one window placed through $[T,2T]$; in the resulting
matrix each distinct zero on the line contributes a positive form of rank one with trace equal to its multiplicity,
and each pair of zeros off the line a block of signature $(1,1)$ \cite[Proposition~4.1, p.~6]{alpoge-furman}, the
latter observation being Bombieri's \cite[§1.3, p.~3]{alpoge-furman}. A rank--trace inequality \cite[Lemma~3.2,
p.~5]{alpoge-furman} and Montgomery's second moment then give the proportions \cite[§1.2, (1.2),
p.~2]{alpoge-furman}. Their matrices are not compressions of an operator with the zeros as spectrum, no trace powers
enter, and their bounds are inequalities for proportions; nothing there states an exact count at a fixed height or
the defect of Theorem~\ref{thm:main}. Lamzouri \cite{lamzouri} gave a second proof of the same proportions with a
single Hilbert-space inequality. Mori \cite{mori} bounds the smallest
eigenvalue of the windowed Weil form from below under the Riemann hypothesis alone, and shows that such bounds force a
proportion of distinct zeros; our defect controls the multiplicity excess itself, not a proportion. Groskin
\cite{groskin} studies finite Galerkin truncations of the Weil form and Kim et al.\ \cite{kim} compute Suzuki's
operator numerically; neither treats multiplicities. On Hexagon, C\'aceres \cite{caceres} states, in the abstract
of his Part~IV, an averaged criterion: the mean energy of the prime-counting error is at least $2B$
unconditionally, with equality if and only if the Riemann hypothesis holds with simple zeros; the test here is
exact at each fixed height, with a certificate, and we have read only that abstract. Numerical verifications of the Riemann hypothesis reach far
greater heights \cite{platt-trudgian}; the certificate here is small, and its interest is that it uses no zero
ordinates.

\paragraph{Scope.}
Theorems~\ref{thm:main} and \ref{thm:integer} are exact equivalences; they give no estimate of $D_{N,k}$ at large
height, so they do not decide whether all zeros are simple. Theorem~\ref{thm:cert} is a finite computation; its
rigour is that of ball arithmetic and of the programs (Section~\ref{sec:cert}), and it is not formally verified.

\section{Proof overview}

The difficulty is to see multiplicity through an operator that cannot see it. In $\Hs=\ell^2(\Gamma,m)$, the operator
$C$ has the single eigenvalue $1/\gamma$ on the coordinate of $\gamma$, whatever $m_\gamma$ is; the weight $m_\gamma$
only scales the norm. So $\tr C^{2k}=2\sum_{\gamma>0}\gamma^{-2k}$, while the Taylor moments $S_k$ of $\log\Xi$
carry the weights $m_\gamma$. Their difference is $2\sum(m_\gamma-1)\gamma^{-2k}\ge0$, which vanishes exactly when
all zeros are simple. The problem is that neither $\Hs$ nor $C$ is finite.

Suzuki's compressions supply the finite objects. Compression to a subspace lowers every singular value, so the finite
defect $D_{N,k}$ is at least the infinite one; nested subspaces make it decrease; and Hilbert--Schmidt convergence,
which holds for any increasing family of subspaces with dense union, makes it converge. Raising $k$ isolates the
largest entry of the sequence $(1/\gamma)$ repeated $m_\gamma-1$ times, which is $1/\gamma_*$. The integer form
counts the eigenvalues of $C_N$ above $1/T$ by Sylvester's law of inertia.

The arithmetic enters only through the matrix entries. For compactly supported test functions the sampling inner
product equals the arithmetic side of the explicit formula, which involves the prime powers below $e^{2N}$ and two
archimedean integrals; the primitive $-i\int f$ realizes $C$ on mean-zero functions. Alp\"oge's theorem is used in one
place: it makes $\Gamma$ real, so that $Q$ is the positive form $\sum m_\gamma|\widehat u(\gamma)|^2$ and $C$ is
self-adjoint. The density of the dyadic windows, which Suzuki's argument needs, is proved directly in
Lemma~\ref{lem:dense}.

\section{The Hilbert space, the windows and the operator}\label{sec:space}

Throughout, the Riemann hypothesis holds, and $\widehat u(t)=\int_\R u(x)e^{itx}dx$. For $u\in L^2(\R)$ with compact
support, $\widehat u$ is entire, and the sampling vector is $Su=(\widehat u(\gamma))_{\gamma\in\Gamma}$.

\begin{lemma}[Sampling bound]\label{lem:sample}
If $u\in H^1(\R)$ is supported in $[-a,a]$, then $|\widehat u(t)|\le\|u'\|_{L^1}/|t|$ for $t\ne0$, and
$\|Su\|_\Hs^2\le 2a\,\sigma\,\|u'\|_{L^2}^2$.
\end{lemma}
\begin{proof}
Integration by parts gives $\widehat{u'}(t)=-it\,\widehat u(t)$, so $|t\widehat u(t)|\le\|u'\|_{L^1}\le
\sqrt{2a}\,\|u'\|_{L^2}$. Summing $m_\gamma|\widehat u(\gamma)|^2\le m_\gamma\gamma^{-2}\cdot2a\|u'\|_2^2$ over
$\Gamma$ gives the second bound.
\end{proof}

\begin{lemma}[Density of the windows]\label{lem:dense}
$V_N\subset V_{N+1}$, and $\bigcup_NSV_N$ is dense in $\Hs$.
\end{lemma}
\begin{proof}
A function linear on the cells of $2^{-N}\mathbb Z$ is linear on the cells of $2^{-N-1}\mathbb Z$, and $(-N,N)\subset
(-N-1,N+1)$; so $V_N\subset V_{N+1}$.

\emph{Test functions are dense.} Let $c\in\Hs$ be orthogonal to $SC_c^\infty(\R)$. The measure
$\tau=\sum_\gamma m_\gamma\overline{c_\gamma}\,\delta_\gamma$ is tempered, since $\sum_\gamma
m_\gamma|c_\gamma|(1+\gamma^2)^{-1}\le\|c\|\sigma^{1/2}<\infty$, and $\tau(\widehat u)=0$ for all $u\in
C_c^\infty(\R)$. The Fourier transform is an automorphism of the Schwartz space $\mathcal S$ and $C_c^\infty$ is dense
in $\mathcal S$, so $\widehat{C_c^\infty}$ is dense in $\mathcal S$. Hence $\tau=0$, and $c=0$.

\emph{The mean can be removed.} Let $\psi\in C_c^\infty(\R)$ with $\int\psi=1$ and $\psi_R(x)=R^{-1}\psi(x/R)$.
Then $\widehat{\psi_R}(t)=\widehat\psi(Rt)$, $|\widehat\psi(t)|\le c_\psi t^{-2}$, and
$\|S\psi_R\|^2=\sum m_\gamma|\widehat\psi(R\gamma)|^2\to0$ as $R\to\infty$ by dominated convergence (each term
tends to $0$ and is at most $c_\psi^2m_\gamma\gamma^{-4}$ for $R\ge1$, which is summable). For $u\in C_c^\infty$,
$u-(\int u)\psi_R$ has mean zero and its sampling vector tends to $Su$.

\emph{Mean-zero test functions are limits of $V_N$.} Let $v\in C_c^\infty$ have mean zero and support in $[-a,a]$, and
let $N\ge a+1$, $h=2^{-N}$. Let $I_Nv$ be the continuous function, linear on the cells of $h\mathbb Z$, that agrees
with $v$ at $h\mathbb Z$. On each cell, $(v-I_Nv)'$ has mean zero, so Wirtinger's inequality gives
$\|(v-I_Nv)'\|_{L^2}\le (h/\pi)\|v''\|_{L^2}$; and, since $I_Nv-v$ vanishes outside $[-a-h,a+h]$,
$|\int I_Nv|=|\int(I_Nv-v)|\le (2a+2h)\,h^2\|v''\|_\infty/12$ by the trapezoid rule. Put $v_N=I_Nv-h^{-1}(\int I_Nv)\,t_0\in V_N$, where $t_0$ is the hat at $0$, of integral $h$ and
with $\|t_0'\|_2^2=2/h$. By Lemma~\ref{lem:sample} with support in $[-a-1,a+1]$,
\[
 \|S(v-v_N)\|\le\sqrt{2(a+1)\sigma}\,\Bigl(\frac h\pi\|v''\|_2+\frac{(2a+2h)h\|v''\|_\infty}{12}\sqrt{2/h}\Bigr)
 \longrightarrow0 .
\]
The three steps give the density.
\end{proof}

On $\Hs$ let $(Cc)_\gamma=c_\gamma/\gamma$. The vectors $e_\gamma=m_\gamma^{-1/2}\mathbf 1_{\{\gamma\}}$ form an
orthonormal basis and $Ce_\gamma=\gamma^{-1}e_\gamma$, so $C$ is bounded and self-adjoint, with
$\|C\|_{S_2}^2=\sum_{\gamma\in\Gamma}\gamma^{-2}\le\sigma$ (as in \cite[§2.3]{suzuki}), and, for $k\ge1$,
\[
 \tr|C|^{2k}=\sum_{\gamma\in\Gamma}\gamma^{-2k}=2\sum_{\gamma>0}\gamma^{-2k}.
\]
Its eigenvalues are $\pm1/\gamma$, one for each distinct ordinate.

\begin{lemma}[The primitive]\label{lem:prim}
For $f\in V_N$, let $F(x)=-i\int_{-\infty}^xf(y)\,dy$. Then $F\in H^1(\R)$ is supported in $[-N,N]$ and $SF=CSf$.
\end{lemma}
\begin{proof}
$F$ vanishes outside $[-N,N]$ because $\int f=0$. From $F'=-if$ and $\widehat{F'}(t)=-it\widehat F(t)$ we get
$\widehat F(t)=\widehat f(t)/t$ for $t\ne0$.
\end{proof}

\section{Proof of Theorems \ref{thm:main} and \ref{thm:integer}}\label{sec:proof1}

\paragraph{The moments.}
$\Xi$ is even, entire of order one, with $\Xi(0)\ne0$ and $\sum m_\gamma\gamma^{-2}<\infty$; its Hadamard product,
paired over $\pm\gamma$, is $\Xi(z)=\Xi(0)\prod_{\gamma>0}(1-z^2/\gamma^2)^{m_\gamma}$ (the exponential factors cancel
in pairs and the remaining factor $e^{bz}$ is $1$ by evenness). For $|z|<\min_{\gamma>0}\gamma$,
$\log(\Xi(z)/\Xi(0))=-\sum_{\gamma>0}m_\gamma\sum_{j\ge1}z^{2j}/(j\gamma^{2j})$, which gives
$S_k=2\sum_{\gamma>0}m_\gamma\gamma^{-2k}$.

\paragraph{The lower bound.}
Write $s_1\ge s_2\ge\cdots$ for singular values. Since $\|P_N\|\le1$, $s_j(P_NCP_N)\le s_j(C)$ for every $j$. The
operator $C_N$ is self-adjoint of finite rank $r$, so
$\tr C_N^{2k}=\sum_{j\le r}s_j(C_N)^{2k}\le\sum_{j\le r}s_j(C)^{2k}<\tr|C|^{2k}$, the last inequality being strict
because $C$ has infinitely many nonzero singular values. Hence
\[
 D_{N,k}>S_k-\tr|C|^{2k}=2\sum_{\gamma>0}\frac{m_\gamma-1}{\gamma^{2k}}\ge\frac{2E(T)}{T^{2k}},
\]
the last step keeping only the ordinates $\gamma\le T$, for which $\gamma^{-2k}\ge T^{-2k}$.

\paragraph{Monotonicity and the limit.}
The subspaces $SV_N$ increase with $N$ (Lemma~\ref{lem:dense}). By the max--min principle, the $n$-th largest
positive eigenvalue of the compression of $C$ to a subspace does not decrease when the subspace grows; applied to $-C$,
the same holds for the absolute values of the negative eigenvalues. So $\tr C_N^{2k}$ is nondecreasing in $N$ and
$D_{N,k}$ is nonincreasing. Since the union of the subspaces is dense, $P_N\to I$ strongly. Following
\cite[(4.1)]{suzuki}: $C-P_NCP_N=(I-P_N)C+P_NC(I-P_N)$, and $\|P_NC(I-P_N)\|_{S_2}\le\|(I-P_N)C\|_{S_2}$ by taking
adjoints; and $\|(I-P_N)C\|_{S_2}^2=\sum_\gamma\gamma^{-2}\|(I-P_N)e_\gamma\|^2\to0$ by dominated convergence. So
$C_N\to C$ in Hilbert--Schmidt norm. For $k\ge1$,
\[
 C^{2k}-C_N^{2k}=\sum_{j=0}^{2k-1}C^{\,j}(C-C_N)C_N^{\,2k-1-j},
\]
and each summand has a factor $C$ or $C_N$ besides $C-C_N$; with $\|UV\|_{S_1}\le\|U\|_{S_2}\|V\|_{S_2}$ and bounded
operator norms, $C_N^{2k}\to C^{2k}$ in trace norm. Hence $D_{N,k}\to S_k-\tr C^{2k}=2\sum(m_\gamma-1)\gamma^{-2k}$.

\paragraph{The supremum.}
Let $v$ be the sequence of the numbers $1/\gamma$, $\gamma>0$, each repeated $m_\gamma-1$ times; it lies in
$\ell^2$, and $\lim_ND_{N,k}=2\|v\|_{2k}^{2k}$. By monotonicity, $\sup_NH_{N,k}=\|v\|_{2k}^{-1}$ (infinite if $v$ is
empty). If $v$ is not empty, its largest entry is $1/\gamma_*$, and
$\gamma_*^{-2k}\le\sum_jv_j^{2k}\le\gamma_*^{-(2k-2)}\sum_jv_j^2$; taking $2k$-th roots, $\|v\|_{2k}\ge1/\gamma_*$ for
every $k$ and $\|v\|_{2k}\to1/\gamma_*$. So $\sup_{N,k}H_{N,k}=\gamma_*$.

\paragraph{The equivalences.}
If $D_{N,k}<2/T^{2k}$, then $E(T)<1$, so $E(T)=0$. Conversely, if $E(T)=0$, then $\gamma_*>T$ (the multiple
ordinates are discrete), and the supremum gives $N,k$ with $H_{N,k}>T$, that is $D_{N,k}<2/T^{2k}$. All zeros are
simple if and only if $\gamma_*=\infty$. This proves Theorem~\ref{thm:main}.

\paragraph{The integer form.}
For $f=\sum_jc_jb_j\in V_N$, $c^*G_Nc=\|Sf\|^2$ and $c^*H_Nc=\langle CSf,Sf\rangle$. The Hermitian form
$c\mapsto c^*(TH_N-G_N)c$ is therefore the pull-back, under the surjection $c\mapsto Sf$ onto $SV_N$, of the form
$x\mapsto\langle(TC_N-I)x,x\rangle$ on $SV_N$. This holds whether or not $G_N$ is invertible: if the map
$c\mapsto Sf$ has a kernel, the form on coefficients vanishes on that kernel, and a subspace on which it is positive
definite meets the kernel only in $0$ and is mapped injectively onto a positive subspace of $SV_N$; conversely a
positive subspace of $SV_N$ lifts to one of the same dimension. So the pull-back has the same positive index. Thus
$\iota_N(T)$ is the number of eigenvalues of $C_N$ on $SV_N$ greater than $1/T$. By max--min these are at most the
eigenvalues of $C$ greater than $1/T$, which are the $1/\gamma$ with $0<\gamma<T$, one per distinct ordinate, and
they do not decrease with $N$. Since $\Xi(T)\ne0$, no eigenvalue of $C$ equals $1/T$; operator-norm convergence
$C_N\to C$ then puts each of the finitely many eigenvalues of $C$ above $1/T$ within reach for large $N$, so
$\iota_N(T)$ eventually equals the number of distinct ordinates in $(0,T]$. That number equals $n(T)$ exactly when
those zeros are simple. This proves Theorem~\ref{thm:integer}.

\begin{proposition}[Size of a certificate]\label{prop:size}
Let $\gamma_1\le\gamma_2\le\cdots$ be the positive ordinates listed with multiplicity. Then
$D_{N,k}\ge2\sum_{j>d_N/2}\gamma_j^{-2k}$. Hence $D_{N,k}<2/T^{2k}$ requires $d_N\ge2n(T)$.
\end{proposition}
\begin{proof}
$\tr C_N^{2k}$ is at most the sum of the $d_N$ largest values of $s_j(C)^{2k}$. These values are $\gamma^{-2k}$, each
distinct $\gamma\in\Gamma$ once, and they are dominated term by term by the list $\pm\gamma_j$ with multiplicity, whose
$d_N$ largest values sum to $2\sum_{j\le d_N/2}\gamma_j^{-2k}$. Subtracting from $S_k$ gives the bound. If
$n(T)>d_N/2$, the term $j=n(T)$ alone contributes at least $2T^{-2k}$.
\end{proof}

\section{The arithmetic evaluation}\label{sec:arith}

\paragraph{The explicit formula.}
For $h\in C_c^\infty(\R)$ define
\begin{align*}
 W(h)={}&\int_\R h(x)\bigl(e^{x/2}+e^{-x/2}\bigr)dx-(\log\pi+\gamma_E)\,h(0)\\
 &+\int_0^\infty\frac{2e^{-2x}h(0)-e^{-x/2}\bigl(h(x)+h(-x)\bigr)}{1-e^{-2x}}\,dx\\
 &-\sum_{n\ge2}\frac{\Lambda(n)}{\sqrt n}\bigl(h(\log n)+h(-\log n)\bigr).
\end{align*}

\begin{lemma}\label{lem:EF}
If $u,v\in H^1(\R)$ have compact support and $h=u*\widetilde v$, $\widetilde v(x)=\overline{v(-x)}$, then
$\langle Su,Sv\rangle_\Hs=W(h)$. Only the prime powers $n$ with $\log n$ in the support of $h$ contribute.
\end{lemma}
\begin{proof}
For $f\in C_c^\infty(0,\infty)$, the explicit formula in Bombieri's form \cite[Appendix ``Explicit formula'', its first three displayed formulas]{cc-weil} reads
$\sum_\rho\widetilde f(\rho)=\int_0^\infty f+\int_0^\infty f^\sharp-\mathcal W_\R(f)-\sum_p\mathcal W_p(f)$, with
$\widetilde f(s)=\int_0^\infty f(x)x^{s-1}dx$, $f^\sharp(x)=x^{-1}f(x^{-1})$,
$\mathcal W_p(f)=\log p\sum_{m\ge1}(f(p^m)+f^\sharp(p^m))$ and
$\mathcal W_\R(f)=(\log4\pi+\gamma_E)f(1)+\int_1^\infty(f(x)+f^\sharp(x)-2x^{-1}f(1))\,dx/(x-x^{-1})$. Take $u,v\in
C_c^\infty$ and $f(e^y)=e^{-y/2}h(y)$. For $\rho=\frac12+i\gamma$, $\widetilde f(\rho)=\widehat h(\gamma)=\widehat
u(\gamma)\overline{\widehat v(\gamma)}$, so the left side is $\langle Su,Sv\rangle$, an absolutely convergent sum by
Lemma~\ref{lem:sample}. With $x=e^y$: $\int f+\int f^\sharp=\int h(y)(e^{y/2}+e^{-y/2})dy$; $f(e^y)+f^\sharp(e^y)=
e^{-y/2}(h(y)+h(-y))$; $dx/(x-x^{-1})=dy/(1-e^{-2y})$; and $\sum_p\mathcal W_p(f)=\sum_n\Lambda(n)n^{-1/2}(h(\log
n)+h(-\log n))$. This gives $W(h)$ with the constant $-(\log4\pi+\gamma_E)$ and $2e^{-y}h(0)$ in the numerator; since
$\int_0^\infty2(e^{-2y}-e^{-y})(1-e^{-2y})^{-1}dy=-\log4$ (substitute $w=e^{-y}$), it equals $W(h)$ as written.

For $u,v\in H^1$ with support in $[-a,a]$, choose $u_n,v_n\in C_c^\infty$ with support in $[-a-1,a+1]$ converging to
$u,v$ in $H^1$. The left sides converge by Lemma~\ref{lem:sample}. The functions $h_n=u_n*\widetilde v_n$ and
$h_n'=u_n'*\widetilde v_n$ converge uniformly to $h$ and $h'$, with supports in a fixed interval. Each term of $W$ is
continuous for this convergence: the first two and the finite prime sum obviously; in the third, on $(0,1]$ the
numerator is at most $c\,x(\|h\|_\infty+\|h'\|_\infty)$ and $1-e^{-2x}\ge(1-e^{-2})x$, and on $[1,\infty)$ the
integrand is at most $c\,e^{-x/2}\|h\|_\infty$.
\end{proof}

\paragraph{The matrices.}
Let $b=t_0*t_0$, the correlation of the central hat:
\[
 b(x)=\begin{cases}2h/3-x^2/h+|x|^3/(2h^2),&|x|\le h,\\ (2h-|x|)^3/(6h^2),&h\le|x|\le2h,\\ 0,&|x|\ge2h,\end{cases}
\]
and $A_d(x)=\int_0^db(x-u)\,du$. Put $g(d)=W(b(\cdot-d))$, which is even in $d$, and $a(d)=W(A_d)$, which is odd
(because $W(\phi(-\,\cdot))=W(\phi)$ and $A_{-d}(x)=-A_d(-x)$). Since $t_j*\widetilde t_i=b(\cdot-(j-i)h)$,
Lemma~\ref{lem:EF} gives
\[
 (G_N)_{ij}=g((j-i)h)-g(jh)-g(ih)+g(0).
\]
For $H_N$, Lemma~\ref{lem:prim} gives $CSb_j=SF_j$ with $F_j=-i\int_{-\infty}^x(t_j-t_0)=i\int_0^{jh}t_0(\cdot-u)\,du$.
Then $F_j*\widetilde t_i=iA_{jh}(\cdot+ih)$, and $A_d(\cdot+c)=A_{d-c}-A_{-c}$, so Lemma~\ref{lem:EF} gives
\[
 (H_N)_{ij}=i\,B_{ij},\qquad B_{ij}=a((j-i)h)-a(jh)+a(ih),
\]
with $B$ real antisymmetric. If $G_N$ is invertible, the matrix of $C_N$ on $SV_N$ in the basis $Sb_j$ is
$G_N^{-1}H_N$, and
\[
 \tr C_N^{2k}=(-1)^k\tr\bigl((G_N^{-1}B)^{2k}\bigr).
\]
Every test function here is supported in $[-2N,2N]$, so only the prime powers $n<e^{2N}$ enter. In the computation,
$A_{mh}=\sum_{0\le j<m}\beta(\cdot-jh)$ with $\beta(x)=\int_0^hb(x-u)\,du$, so that each functional is evaluated on
a piecewise polynomial supported on five cells.

\section{The certificate}\label{sec:cert}

Table~\ref{tab:cert} gives, for six windows, an enclosure of $D_{N,k}$ and the threshold $2/T^{2k}$. Each row
satisfies the strict test $D_{N,k}<2/T^{2k}$, so by Theorem~\ref{thm:main} every zero with $0<\Im\rho\le T$ is simple;
the last row is Theorem~\ref{thm:cert}. Each height is $H_{N,k}$ rounded down, and $d_N\ge2n(T)$ as
Proposition~\ref{prop:size} requires ($n(97.632)=28$).

\begin{table}[ht]
\centering\small
\begin{tabular}{@{}rrrrrll@{}}
\toprule
$N$&$d_N$&prime powers&$k$&$T$&$D_{N,k}$ (enclosure)&$2/T^{2k}$\\\midrule
1&2&5&9&14.134&$3.94653014372982665\ldots\cdot10^{-21}$&$3.9469209\ldots\cdot10^{-21}$\\
2&14&24&3&20.704&$2.53910511009839593\ldots\cdot10^{-8}$&$2.5392440\ldots\cdot10^{-8}$\\
3&46&98&3&26.782&$5.41964807266176611\ldots\cdot10^{-9}$&$5.4196589\ldots\cdot10^{-9}$\\
4&126&465&2&39.913&$7.88033023780063053\ldots\cdot10^{-7}$&$7.8808399\ldots\cdot10^{-7}$\\
5&318&2532&2&62.413&$1.31802357404210376\ldots\cdot10^{-7}$&$1.3180435\ldots\cdot10^{-7}$\\
6&766&15040&3&97.632&$2.30927091609723675\ldots\cdot10^{-12}$&$2.3092815\ldots\cdot10^{-12}$\\\bottomrule
\end{tabular}
\caption{Strict height certificates. The enclosure radii are below $10^{-26}$ in every row; the digits shown are
inside the enclosure.}\label{tab:cert}
\end{table}

\paragraph{How the enclosures were computed.}
All quantities are evaluated in Arb ball arithmetic \cite{arb}, which returns for each real number an interval
guaranteed to contain it. (i) The functionals $g(d)$ and $a(d)$: on each cell of $h\mathbb Z$ the test function is a
polynomial of degree at most four, recovered exactly by interpolation; the first, third and archimedean integrals are
evaluated by rigorous ball-arithmetic integration \cite{arb-int}, the integrand on $[0,h]$ being written as
$\bigl(-3h(0)e^{-x/2}\varphi(-3x/2)-e^{-x/2}\psi_1(x)\bigr)/\bigl(2\varphi(-2x)\bigr)$ with $\varphi(z)=(e^z-1)/z$
and $h(x)+h(-x)=2h(0)+x\psi_1(x)$, which removes the singularity at $0$; the tail $\int_h^\infty
2e^{-2x}(1-e^{-2x})^{-1}dx=-\log(1-e^{-2h})$ is exact; the prime sum is finite. (ii) $G_N$ and $B$ are assembled
from these balls; a rigorous linear solve certifies that $G_N$ is invertible (so it is positive definite, being a
Gram matrix) and encloses $G_N^{-1}B$; the trace of its $2k$-th power is enclosed. (iii) $S_k$ is the coefficient of
$z^{2k}$ in the power series of $\log(s(s-1)/2)-\tfrac s2\log\pi+\log\Gamma(s/2)+\log\zeta(s)$ at $s=\frac12+iz$,
computed in ball arithmetic. No zero of $\zeta$ is used.

The computation was done twice, by two independently written programs: once by the producer of the mathematics,
at 120 digits with a preconditioned positive-definiteness test, and once for this note, at 150 digits, by the program
on the workbench \cite{wb}. The enclosures agree. The rigour of the certificate is that of ball arithmetic and of
these programs; it is not formally verified.

\section{Explicit constants}

\begin{proof}[Proof of Proposition~\ref{prop:const}]
Let $\gamma\ge10$ be an ordinate of multiplicity $m$ and $s=2+i\gamma$. Pairing the Hadamard product of $\xi$,
$\Re\frac{\xi'}{\xi}(s)=\sum_\rho\Re\frac1{s-\rho}$, and every term is positive because $\Re\rho=\frac12$; the $m$
terms with $\rho=\frac12+i\gamma$ give $m/(3/2)$. So $\frac{2m}3\le\Re\frac{\xi'}{\xi}(2+i\gamma)$. Now
$\frac{\xi'}{\xi}(s)=\frac1s+\frac1{s-1}-\frac12\log\pi+\frac12\frac{\Gamma'}{\Gamma}(\frac s2)+\frac{\zeta'}{\zeta}(s)$.
Here $\Re(\frac1s+\frac1{s-1})<3/\gamma^2$; $|\frac{\zeta'}{\zeta}(s)|\le\sum\Lambda(n)n^{-2}<0.57687<\frac35$
(summing the prime powers up to $1000$ and bounding the rest by $\int_{1000}^\infty(\log x)x^{-2}dx=
(\log1000+1)/1000$); and with $y=\gamma/2$, $\Re\frac{\Gamma'}{\Gamma}(1+iy)=-\gamma_E+\sum_{n\ge1}\frac{y^2}{n(n^2+y^2)}
\le-\gamma_E+\frac{y^2}{1+y^2}+\frac12\log(1+y^2)$, comparing the decreasing summand with its integral. Collecting,
\[
 \Re\frac{\xi'}{\xi}(2+i\gamma)<\frac12\log\gamma+\frac4{\gamma^2}+\frac{11}{10}-\frac12(\gamma_E+\log2\pi)
 <\frac12\log\gamma+\frac1{25}-\frac1{20}<\frac12\log\gamma,
\]
using $\frac14\log(1+4/\gamma^2)<1/\gamma^2$, $\gamma_E>\sum_{j\le6}j^{-1}-\log7>\frac12$ and $\log2\pi>\frac95$.
Hence $m<\frac34\log\gamma$, that is $\gamma>e^{4m/3}$. Theorem~\ref{thm:cert} gives $\gamma_*>97.632$, and
$97.632^2/2=4766.003712$.
\end{proof}

\section{The backward heat flow}\label{sec:heat}

Let $\Phi(u)=\sum_{n\ge1}(2\pi^2n^4e^{9u}-3\pi n^2e^{5u})e^{-\pi n^2e^{4u}}$ and, for $t\in\R$,
$H_t(z)=\int_0^\infty e^{tu^2}\Phi(u)\cos(zu)\,du$; then $H_0(z)=\frac18\xi(\frac12+\frac{iz}2)=\frac18\Xi(z/2)$
\cite[§1, (1)--(4)]{rodgers-tao}, a zero of $\Xi$ at $\gamma$ is a zero of $H_0$ at $x=2\gamma$, and differentiation under
the integral gives $\partial_tH_t=-\partial_z^2H_t$. For $a>0$ write
\[
 F_a(w)=\frac{H_{-a}(i\sqrt w)}{H_{-a}(0)}=\sum_{k\ge0}b_k(a)w^k,\qquad
 J_{d,a}(X)=\sum_{k=0}^d(d)_k\,b_k(a)X^k,
\]
with the falling factorial $(d)_k=d(d-1)\cdots(d-k+1)$, and let $j_*(-a)\in[1,\infty]$ be the least $d$ for which
$J_{d,a}$ has a nonreal zero ($\infty$ if there is none).

\begin{proposition}\label{prop:heat}
Assume the Riemann hypothesis. Then $\lim_{a\downarrow0}a\,j_*(-a)=\gamma_*^2/2$, the limit being $+\infty$ when all
zeros are simple. With Theorem~\ref{thm:main} this reads
$\lim_{a\downarrow0}a\,j_*(-a)=\frac12(\sup_{N,k}H_{N,k})^2$, and Proposition~\ref{prop:const} gives the value
$>4766.003712$.
\end{proposition}

The proof occupies the rest of this section. Write $x_*=2\gamma_*$, so that $\gamma_*^2/2=x_*^2/8$, and let
$m_x$ be the multiplicity of the zero $x$ of $H_0$.

\paragraph{Growth and products.}
Each term of $\Phi$ is positive for $u\ge0$, since $2\pi n^2e^{4u}\ge2\pi>3$; so $H_{-a}(iy)>0$ for real $y$, $b_k(a)>0$,
and $F_a>0$ on $[0,\infty)$. Splitting the integral at $u=\log(R+3)$, beyond which $e^{-\pi e^{4u}}$ dominates,
gives $\log\max_{|z|\le R}|H_t(z)|=O(R\log(R+3))$ uniformly for $t$ in compact sets. Hence $G(a,w)=H_{-a}(i\sqrt w)$ is
entire in $(a,w)$ (it is even in $\sqrt w$), of order at most $\frac12$ in $w$, and $F_a(w)=\prod_i(1-w/w_i)$ over its
zeros, which are $w_i=-z_i^2$ for the zeros $z_i$ of $H_{-a}$ with $\Re z_i>0$; none lies on $[0,\infty)$. At $a=0$,
by the Riemann hypothesis, $H_0(z)=H_0(0)\prod_{x>0}(1-z^2/x^2)^{m_x}$ and, by Jensen's formula and the growth bound,
$N_H(R)=\sum_{0<x\le R}m_x\ll R\log(R+3)$.

\paragraph{The zeros of $H_{-a}$ lie in thin cones.}
For nonreal $z$ put $L(z)=H_0'(z)/H_0(z)$ and $S(z)=\sum_{x>0}m_x(|z-x|^{-2}+|z+x|^{-2})$. The entire function
$Q_z(v)=H_0(z+v)/(H_0(z)e^{L(z)v})$ is the product of $(1+v/(z-\xi))e^{-v/(z-\xi)}$ over the zeros $\xi=\pm x$, so
$|1+w|^2\le e^{2\Re w+|w|^2}$ gives $|Q_z(v)|\le e^{S|v|^2/2}$; with $Q_z(0)=1$, $Q_z'(0)=0$ and Cauchy's estimate this gives
$|Q_z(v)-1|\le8S|v|^2e^{2S|v|^2}$. Since $H_{-a}$ is the convolution of $H_0$ with the Gaussian of variance $2a$, moving the
line of integration by $p=2aL(z)$ and completing the square gives
$H_{-a}(z)/(H_0(z)e^{aL(z)^2})=\mathbb EQ_z(p+U)$ with $U\sim N(0,2a)$. With $A=aS(z)$, $B=a^2|L(z)|^2S(z)$ and
$|p+U|^2\le2|p|^2+2U^2$, the Gaussian moments give, for $16A<1$,
\[
 \Bigl|\frac{H_{-a}(z)}{H_0(z)e^{aL(z)^2}}-1\Bigr|\le16e^{16B}\Bigl(\frac{4B}{\sqrt{1-16A}}+\frac{2A}{(1-16A)^{3/2}}\Bigr),
\]
which is at most $\frac7{16}$ when $A,B\le2^{-10}$; then $H_{-a}(z)\ne0$. If $|z|=r$ and $|\Im z|\ge q r$, splitting the
zeros at $r/2$ and $2r$ and using $N_H$ gives $S(z)\ll q^{-2}\log(r+3)/(r+1)$ and $|L(z)|\ll q^{-1}\log(r+3)$. Fix
$\eta\in(0,\frac12)$ and take $q=\eta\sqrt a$: then $A\ll\eta^{-2}\log(r+3)/(r+1)$ and
$B\ll\eta^{-4}\log^3(r+3)/(r+1)$, uniformly in $0<a\le1$. Hence, for every $\kappa>0$ there is $R_0$ such that every
zero of $H_{-a}$, $0<a\le1$, with $|z|\ge R_0$ and $\Re z>0$ has $|\arg z|\le\kappa\sqrt a$.

\paragraph{Jensen polynomials as differential operators.}
For a power series $F=\sum f_kw^k$, $F(D)z^d=\sum_kf_k(d)_kz^{d-k}=z^dJ_d(1/z)$ with $J_d(X)=\sum_{k\le d}(d)_kf_kX^k$; so
$J_d$ has only real zeros if and only if $F(D)z^d$ does. By the theorem of Borcea and Br\"and\'en
\cite[Theorem~2(b), p.~5]{bb}, a real linear operator on polynomials of degree at most $d$ preserves real-rootedness if
its symbol $T[(z+w)^d]$ has no zero with $\Im z,\Im w>0$; this is their condition~(b), since $H_2(\R)$ there
denotes the real stable polynomials in two variables, those with real coefficients and no zero with both
imaginary parts positive. For $T=1+\beta D$ with $\beta>0$ the symbol is
$(z+w)^{d-1}(z+w+d\beta)$. For $T=1+2\beta\cos\theta\,D+\beta^2D^2$, $|\theta|<\pi/2$, it is
$(z+w)^{d-2}\bigl((z+w)^2+2d\beta\cos\theta(z+w)+d(d-1)\beta^2\bigr)$, whose quadratic factor in $z+w$ has discriminant
$4d\beta^2(1-d\sin^2\theta)$ and hence real roots when $d\sin^2\theta\le1$ (for $d=1$ the symbol is
$z+w+2\beta\cos\theta$). Since $\Im(z+w)>0$, these symbols do not vanish. Compositions of such operators preserve
real-rootedness in degree at most $d$, and so do coefficientwise limits, by Hurwitz's theorem, when the limit is not
zero.

\paragraph{The interpolation and its local form.}
Let $G(a,w)=\sum_kc_k(a)w^k$ and $\mathcal J(a,\epsilon,w)=\sum_kc_k(a)w^k\prod_{j<k}(1-j\epsilon)$. The growth bound and
Cauchy's estimate give $|c_k(a)|\le C^kk^{-4k/3}$ for $a$ in compact sets, and the product is at most $(1+E)^kk!$ for
$|\epsilon|\le E$; so $\mathcal J$ is entire in $(a,\epsilon,w)$, and $\mathcal J(a,1/d,w)=H_{-a}(0)J_{d,a}(w/d)$. Expanding
$\prod_{j<k}(1-j\epsilon)=\sum_s(-\epsilon)^se_s(k)$, where $e_s(k)$, the elementary symmetric function of
$0,\ldots,k-1$, is a polynomial in $k$ of degree $2s$ with leading coefficient $1/(2^ss!)$ (by
$e_s(k+1)-e_s(k)=ke_{s-1}(k)$), and writing it in falling factorials of $k$, the coefficient of $\epsilon^s$ is a
differential operator in $w$ of order at most $2s$ with principal part $(-1)^sw^{2s}D^{2s}/(2^ss!)$. In the variable $w$,
$\partial_aG=\mathcal LG$ with $\mathcal L=-4wD^2-2D$, so the coefficient of $a^r$ in $G$ is $\mathcal L^rG(0,\cdot)/r!$,
with principal part $(-4w)^rD^{2r}/r!$.

\emph{Local form.} Let $K(a,w)=\sum_kK_k(a)w^k$ be holomorphic for $a$ near $0$ and $w$ in a disk about $0$ that
contains $w_0$, and suppose that its interpolation $\mathcal J_K(a,\epsilon,w)=\sum_kK_k(a)w^k\prod_{j<k}(1-j\epsilon)$
is holomorphic in $(a,\epsilon,w)$ there and entire in $\epsilon$ (for $G$ this was shown above; for a polynomial it is
clear). Suppose that, for each $r$, the coefficient of $a^r$ in $K$
vanishes at $w_0$ to order at least $m-2r$, with coefficient of $(w-w_0)^{m-2r}$ equal to
$A_0(-4w_0)^r\,m!/(r!(m-2r)!)$, $A_0\ne0$. Apply the interpolation $\mathcal J_K$ (the same multipliers applied to the
Taylor coefficients of $K$ in $w$) and substitute $a=\gamma q^2$, $\epsilon=q^2$, $w=w_0+qu$. A term
$a^r\epsilon^s$ lowers the vanishing order by at most $2s$ and carries $q^{2r+2s}$, so no power below $q^m$ occurs, and
at $q^m$ only the principal parts contribute:
\[
 \frac{\mathcal J_K(\gamma q^2,q^2,w_0+qu)}{A_0q^m}=\sum_{k}\ \sum_{r+s=k}\frac{(-4\gamma w_0)^r}{r!}\frac{(-w_0^2/2)^s}{s!}
 D_u^{2k}u^m+O(q)=e^{c D_u^2}u^m+O(q),
\]
with $c=4\gamma x^2-x^4/2$ when $w_0=-x^2$, uniformly for $\gamma$ and $u$ in compact sets. The function $K=G$ satisfies
the hypothesis at $w_0=-x^2$ with $m=m_x$, by the principal parts of $\mathcal L^r$. If $c>0$, then
$e^{cD^2}u^m=(i\sqrt{2c})^m\mathrm{He}_m(u/(i\sqrt{2c}))$ has the nonreal roots $i\sqrt{2c}\,h$ for the nonzero roots
$h$ of the Hermite polynomial $\mathrm{He}_m$, which exist for $m\ge2$; if $c<0$, it is
$(\sqrt{2|c|})^m\mathrm{He}_m(u/\sqrt{2|c|})$, with $m$ simple real roots.

\paragraph{Upper bound.}
Let $x_*<\infty$ and $\gamma>x_*^2/8$, so $c>0$ at $x=x_*$, where $m\ge2$. Put $d=\lceil\gamma/a\rceil$ and $q=d^{-1/2}$; then
$a=(ad)q^2$ with $ad\to\gamma$. By the local form and Rouch\'e's theorem, $\mathcal J(a,1/d,\cdot)$, and hence
$J_{d,a}$, has a nonreal zero for all small $a$. So $\limsup_{a\downarrow0}a\,j_*(-a)\le\gamma$ for every
$\gamma>x_*^2/8$.

\paragraph{Lower bound.}
Fix $0<\gamma_0<x_*^2/8$ (any $\gamma_0>0$ if $x_*=\infty$), put $\kappa=1/(4\sqrt{\gamma_0})$ and take $R_0$ from the
cone estimate. Choose $R>R_0^2$ with no zero of $F_0$ on $|w|=R$; for small complex $a$ the circle encloses a fixed
number $\delta$ of zeros of $F_a$, the power sums $\frac1{2\pi i}\oint w^kF_a'/F_a\,dw$ are holomorphic in $a$, and Newton's
identities give a polynomial $P_a$ with these zeros and $P_a(0)=1$, real for real $a$. Then $F_a=P_aT_a$ with
$T_a=\prod(1-w/w_i)$ over the zeros outside the circle. Near an enclosed endpoint $w_0=-x^2$, $P_a=UG$ with $U=1/(G(a,0)T_a)$
holomorphic and nonvanishing; multiplying by $U=\sum_ja^jU_j$ only raises vanishing orders in the terms with
$j\ge1$, so $P_a$ satisfies the hypothesis of the local form with $A_0=AU_0(w_0)$.

\emph{The head.} Apply the local form to $P_a$ with $\gamma\in[0,\gamma_0]$. At a multiple endpoint $x\ge x_*$, so
$c=4\gamma x^2-x^4/2<0$; at a simple endpoint the limit is $u$. Rouch\'e disks about these $\delta$ real roots, uniform in
$\gamma$, each contain one root, which is real because the coefficients are real; they exhaust the $\delta$ roots of
the degree-$\delta$ polynomial. So there is $d_1$ with: $d\ge d_1$ and $ad\le\gamma_0$ imply that $J_d(P_a)$ has only real
zeros. For $d<d_1$: $P_0=\prod_i(1+\beta_iw)$ with $\beta_i>0$, and applying $1+\beta D$ to a polynomial whose nonzero roots
are simple and negative gives one of the same kind (the equation $p'/p=-1/\beta$ has one root in each gap and one to
the left, $p'/p$ being decreasing on each interval). So each $J_d(P_0)$, $d<d_1$, has simple real roots, which persist
for small real $a$.

\emph{The tail.} The zeros of $T_a$ are $w=-z^2$ with $|z|>R_0$; real $z$ give factors $1+\beta w$ with $\beta>0$, and
the conjugate pairs give factors $1+2\beta\cos\theta\,w+\beta^2w^2$ with $|\theta|=2|\arg z|\le2\kappa\sqrt a$. For
$d\le\gamma_0/a$, $d\sin^2\theta\le4\kappa^2\gamma_0=\frac14$. By the operator criterion, $T_a(D)$ preserves real-rootedness
in degree at most $d$; applied to $P_a(D)z^d$, which has only real zeros by the head, it shows that $F_a(D)z^d$, and so
$J_{d,a}$, has only real zeros for every $d\le\gamma_0/a$ and all small $a$. Hence
$\liminf_{a\downarrow0}a\,j_*(-a)\ge\gamma_0$. Letting $\gamma_0\uparrow x_*^2/8$, or $\gamma_0\to\infty$ when
$x_*=\infty$, and combining with the upper bound proves Proposition~\ref{prop:heat}.\qed

\section{What remains open}

All zeros are simple if and only if $\sup_{N,k}H_{N,k}=\infty$, equivalently if and only if, at an unbounded sequence
of heights $T$ with $\Xi(T)\ne0$, the integer deficit $n(T)-\ind_+(TH_N-G_N)$ vanishes for some $N$. A proof would need
an upper estimate for $D_{N,k}$ at unbounded heights, with the dimension $d_N\ge2n(T)$ that
Proposition~\ref{prop:size} forces. No such estimate is known to us.

\section*{Note on method}
The mathematics of Theorems~\ref{thm:main}--\ref{thm:cert} and Propositions~\ref{prop:const}, \ref{prop:size}
and~\ref{prop:heat} was
produced by GPT-6 Astra (ChatGPT Pro mode), with the first computation of the certificate. Claude Opus~5.5 at effort
high checked every proof and wrote this text, including the proofs of Lemma~\ref{lem:dense} and
Lemma~\ref{lem:EF}, and wrote the second program for the certificate; a third, floating-point computation of all six
rows by the same model agrees with both. A separate instance of Claude Opus~5.5 at effort high, given only the
statements of Theorems~\ref{thm:main} and~\ref{thm:integer} and Suzuki's paper, re-proved both independently and
then compared its proofs with ours; it found no gap, and the proof that the $b_j$ form a basis is taken from it.
Suzuki's paper was read at the
statements cited, Connes and Consani's appendix at the explicit formula, and the abstracts of
\cite{alpoge-furman,lamzouri,mori,groskin,kim,platt-trudgian}; the theorem of Borcea and Br\"and\'en at its statement, and Rodgers and Tao
at the definitions of the heat family. The same separate instance of Claude Opus~5.5 at effort high then checked
the proof of Proposition~\ref{prop:heat} step by step, opening both of these sources at their locators, and found no
error. No human has checked this note. Our text and programs are dedicated to the public domain under CC0~1.0;
the permanent record of this work is WB-0009, \expandafter\url\expandafter{\recordurl}; its programs and check
records are at \expandafter\url\expandafter{\workbenchurl} \cite{wb}.

\begin{thebibliography}{99}\small
\bibitem{alpoge} L.~Alp\"oge, announcement of the proof of the Riemann hypothesis and of the generalized Riemann
hypothesis for Dirichlet $L$-functions, 20 November 2026, \url{https://alpo.ge/remark.pdf}.
\bibitem{alpoge-furman} L.~Alp\"oge and R.~Furman, More than two thirds of the zeta zeros are simple and on the
critical line, arXiv:2608.13637v2 (2026); page numbers refer to that version.
\bibitem{arb} F.~Johansson, Arb: efficient arbitrary-precision midpoint-radius interval arithmetic, \emph{IEEE Trans.
Comput.} 66 (2017), 1281--1292.
\bibitem{arb-int} F.~Johansson, Numerical integration in arbitrary-precision ball arithmetic, arXiv:1802.07942
(2018).
\bibitem{bb} J.~Borcea and P.~Br\"and\'en, P\'olya--Schur master theorems for circular domains and their
boundaries, \emph{Ann. of Math.} 170 (2009), 465--492; arXiv:math/0607416v6, whose page numbers are used.
\bibitem{caceres} P.~C\'aceres, The Riemann hypothesis as a margin: a two-sided criterion verified to $2\cdot10^{29}$,
the fourth-moment anatomy of M\"obius cancellation, and one constant seen from every side, hexagon:2610.00688v3
(2026), abstract, \url{https://hexagonmath.org/2610.00688}.
\bibitem{cc-weil} A.~Connes and C.~Consani, Weil positivity and trace formula, the archimedean place, Appendix
``Explicit formula'', following E.~Bombieri, The Riemann hypothesis, in \emph{The Millennium Prize Problems}, Clay
Math. Inst., 2006, 107--124.
\bibitem{groskin} A.~Groskin, A finite Guinand--Weil dictionary and archimedean tail order for the truncated Weil
quadratic form, arXiv:2607.02828v3 (2026).
\bibitem{kim} T.~Kim, Y.~Hong, M.~Kim, S.~Choi, J.~Jang and M.~Kim, A numerical realization of Suzuki's
Weil-quadratic-form operator, arXiv:2607.24830v3 (2026).
\bibitem{lamzouri} Y.~Lamzouri, A new proof that more than $2/3$ of the zeros of the Riemann zeta function are simple
and on the critical line, arXiv:2609.02882v2 (2026).
\bibitem{lean} Lean FRO and Mathlib Initiative, Lean verification of the Riemann Hypothesis, 21 November 2026,
\url{https://leanprover-community.github.io/blog/posts/riemann-hypothesis/}.
\bibitem{mori} R.~Mori, Windowed Weil forms under RH alone: an $e^{-\mu^{6/5+\varepsilon}}$ lower bound and a barrier
from distinct zeros, Zenodo 24281606, version 1 (2026).
\bibitem{platt-trudgian} D.~Platt and T.~Trudgian, The Riemann hypothesis is true up to $3\cdot10^{12}$, \emph{Bull.
London Math. Soc.} 53 (2021), 792--797.
\bibitem{rodgers-tao} B.~Rodgers and T.~Tao, The de Bruijn--Newman constant is non-negative, \emph{Forum Math. Pi} 8
(2020), e6; arXiv:1801.05914v5, §1.
\bibitem{suzuki} M.~Suzuki, Spectral convergence associated with Weil's quadratic form, arXiv:2610.32704 (2026).
\bibitem{wb} WB-0009. Permanent record:\newline \expandafter\url\expandafter{\recordurl};\newline programs and
check records:\newline \expandafter\url\expandafter{\workbenchurl}.
\end{thebibliography}
\end{document}
